% =====================================================================
% CODEX ColombiaMacro · Parte I · Cómo usar ColombiaMacro (español)
% Fragmento para \input: sin preámbulo.
% =====================================================================
\providecommand{\cmsym}[1]{{\fontspec{DejaVu Sans}#1}}
\providecommand{\cmexp}{{\fontspec{DejaVu Serif}⤢}}
\newcolumntype{L}[1]{>{\raggedright\arraybackslash}p{#1}}

\part{Cómo usar ColombiaMacro}

% ---------------------------------------------------------------------
\chapter{La pregunta que responde el sitio}
\label{es:u:pregunta}

ColombiaMacro es un tablero académico e institucional sobre la economía colombiana, construido
exclusivamente con datos oficiales. Todo el sitio está organizado alrededor de una sola pregunta:

\begin{center}
\sffamily\large ¿En qué fase del ciclo está la economía colombiana\\
y qué implica para los precios, la política monetaria y los activos?
\end{center}

La pregunta se descompone en preguntas encadenadas, una por página de análisis. El orden no es casual:
sigue la secuencia con la que un banco central y un inversionista institucional leen una economía:
\emph{actividad} (¿cuánto produce el país y frente a qué capacidad?) $\rightarrow$ \emph{precios}
(¿qué presión genera esa actividad sobre la inflación?) $\rightarrow$ \emph{reacción de política}
(¿qué hace el Banco de la República con su tasa?) $\rightarrow$ \emph{precios de los activos} (curva de
los TES, tasa de cambio, bolsa) $\rightarrow$ \emph{restricciones} (cuentas externas, cuentas fiscales y
comercio exterior).

\section{Principios con los que se construye cada cifra}

\begin{enumerate}
\item \textbf{Solo datos oficiales.} DANE, Banco de la República (BanRep), Ministerio de Hacienda,
  Superintendencia Financiera, Superintendencia de Sociedades, DIAN (vía DANE), BVC/MSCI (vía BanRep) y,
  para la composición de la bolsa, el fondo que replica el índice COLCAP. Ninguna cifra se digita a mano.
\item \textbf{Sin proyecciones propias.} El sitio describe lo observado. Cuando muestra algo que mira hacia
  adelante (por ejemplo, la inflación que descuentan los bonos del Gobierno), es una lectura \emph{del
  mercado}, calculada con precios observados, y se identifica como tal.
\item \textbf{Frases con respaldo.} Cada frase automática (resúmenes, etiquetas de estado, el veredicto) sale
  de una regla explícita con umbrales fijos, y cita el dato que la sustenta. Las reglas están en el
  capítulo~\ref{es:i:estados}.
\item \textbf{Un eje, una unidad.} Cada gráfico usa una sola escala vertical, salvo excepciones rotuladas en
  las que una misma magnitud se muestra en otra unidad.
\item \textbf{Incertidumbre visible.} Lo que se estima (la capacidad de la economía, la tasa neutral, la
  inflación implícita) se presenta con su rango o con su advertencia, no como una cifra exacta.
\item \textbf{Sin información futura.} Los cálculos ``en tiempo real'' solo usan la información que existía en
  cada fecha.
\end{enumerate}

\section{Para quién es}
\begin{itemize}
\item \textbf{Inversionistas extranjeros}: una lectura en inglés o en español con las variables que importan
  para una asignación en Colombia (tasas reales, expectativas de inflación, tasa de cambio, bolsa en dólares,
  cuentas externas y fiscales).
\item \textbf{Inversionistas, analistas y empresarios locales}: un tablero con la disciplina metodológica de
  un informe de banco central.
\item \textbf{Comunidad académica}: todas las fórmulas son explícitas (Parte~II) y cada gráfico declara su
  fuente y su método.
\end{itemize}

\begin{alerta}[Lo que el sitio no es]
ColombiaMacro no es una recomendación de inversión ni un pronóstico. Las etiquetas (``Alta'', ``Peso más
fuerte'', ``Expansión''\ldots) \emph{describen} cada dato frente a una regla; no dicen qué hacer. Las
magnitudes estimadas (brecha del producto, tasa neutral, inflación implícita) tienen incertidumbre
documentada en el capítulo~\ref{es:r:limitaciones}.
\end{alerta}

% ---------------------------------------------------------------------
\chapter{Mapa del sitio}
\label{es:u:mapa}

El sitio tiene una \textbf{portada} (el tablero), \textbf{once páginas de análisis} agrupadas en tres
bloques del menú y una página de \textbf{Datos}. La versión en español vive en la raíz del sitio y la versión
en inglés en \texttt{/en/}; ambas tienen exactamente la misma estructura.

\section{El menú}
La barra superior muestra el logotipo (que lleva a la portada), tres menús desplegables y el enlace a Datos:

\begin{center}\small
\begin{tabularx}{\linewidth}{@{}lX@{}}
\toprule
\textbf{Grupo del menú} & \textbf{Páginas}\\\midrule
Actividad & Ciclo · Crecimiento · Capacidad · Empleo\\
Precios y tasas & Inflación · Tasas · Curva TES\\
Mercados y cuentas & Tasa de cambio · Sector empresarial · Externo y fiscal · Comercio exterior\\
Datos & Todos los indicadores · publicaciones oficiales · fuentes · descargas CSV\\
\bottomrule
\end{tabularx}
\end{center}

A la derecha están el enlace de idioma (\emph{English}/\emph{Español}), el botón de tema claro u oscuro y, en
pantallas pequeñas, el botón \cmsym{☰} \emph{Secciones}, que abre el menú. Cada página de análisis lleva
además una ruta de navegación (\emph{Tablero / Actividad / Ciclo}) y, al final, los enlaces
\emph{Anterior}, \emph{Volver al tablero} y \emph{Siguiente}, que recorren las páginas en el orden del menú.

\section{Las once páginas de análisis y la pregunta de cada una}

\begin{longtable}{@{}L{0.19\linewidth}L{0.33\linewidth}L{0.40\linewidth}@{}}
\toprule
\textbf{Página} & \textbf{Pregunta central} & \textbf{Qué se encuentra}\\\midrule\endhead
\textbf{Ciclo}\newline\small(\emph{El ciclo económico}) &
¿En qué fase del ciclo está la economía, cómo llegó ahí y hacia dónde se mueve? &
Fase trimestral y mensual; consenso de cinco métodos de brecha del producto; reloj del ciclo interactivo;
amplitud sectorial (cuántos sectores están sobre su tendencia); fechado de expansiones y recesiones; ciclo
y empleo (ley de Okun).\\
\textbf{Crecimiento}\newline\small(\emph{Crecimiento, sectores y demanda}) &
¿Está creciendo la economía, a qué velocidad y por qué? &
Tres velocidades del PIB (anual, 12 meses, trimestral anualizada); crecimiento real frente a precios
(deflactor); aportes por sector y por componente de la demanda; los 12 sectores trimestre a trimestre;
inversión; consumo de los hogares; PIB por persona y productividad; nivel frente al camino previo a 2020.\\
\textbf{Capacidad}\newline\small(\emph{Capacidad productiva}) &
¿Produce la economía por encima o por debajo de lo que puede sostener sin generar inflación, y dónde? &
PIB frente a su capacidad (PIB potencial); brecha por sector; holgura del mercado laboral (desempleo,
subempleo, subutilización); regiones: crecimiento, tamaño frente a 2019, estructura productiva y
diversificación de cada departamento.\\
\textbf{Empleo}\newline\small(\emph{Empleo e informalidad}) &
¿Cuántos trabajan, dónde y cómo? &
Desempleo e informalidad (nacional, 13 y 23 ciudades, por sector); empleos creados o perdidos por rama;
tipo de empleo (posición ocupacional); mujeres y hombres; jóvenes; campo y ciudad; población fuera de la
fuerza de trabajo; desempleo en las 32 capitales.\\
\textbf{Inflación}\newline\small(\emph{Inflación y precios}) &
¿Cuánto suben los precios frente a la meta, qué espera el mercado y de dónde viene la inflación? &
Inflación total y de fondo frente a la meta; grupos de precios; inflación esperada a 1 año y trayectoria a
10 años (5y5y); aportes por división; bienes frente a servicios; distribución entre 188 subclases;
inflación por nivel de ingreso y en 23 ciudades.\\
\textbf{Tasas}\newline\small(\emph{Tasas de interés}) &
¿Qué está haciendo el Banco de la República y cuánto cuesta de verdad endeudarse? &
Tasa de política frente a inflación; tasa real ex ante frente a la neutral; ciclos de la tasa desde 2000;
traspaso a tasas de ahorro y crédito; costo del crédito por modalidad; mercado monetario (IBR);
crecimiento real del crédito; liquidez del Banco.\\
\textbf{Curva TES}\newline\small(\emph{Curva de rendimientos}) &
¿Cuánto cobra el mercado por prestarle al Gobierno a 1, 5 y 10 años? &
Explorador de la curva cero cupón (pesos y UVR) para cualquier día desde 2003; nivel, pendiente y
curvatura; episodios de curva invertida; cambios por plazo; descomposición en tasa real y compensación por
inflación; prima sobre la tasa del Banco; volatilidad.\\
\textbf{Tasa de cambio}\newline\small(\emph{El peso colombiano}) &
¿Cómo está el peso, es el peso o es el dólar, y qué lo mueve? &
TRM; tasa de cambio real (ITCR); el peso frente al dólar global y a monedas vecinas; petróleo y términos de
intercambio; balanza cambiaria y reservas; volatilidad.\\
\textbf{Sector empresarial}\newline\small(\emph{Las empresas de Colombia}) &
¿Qué empresas mueven la bolsa y cómo están las empresas grandes? &
COLCAP frente a un índice equiponderado y a las 7 Magníficas; peso de cada acción y sector; las 10.000
empresas más grandes (ingresos, utilidades, rentabilidad, endeudamiento, concentración, sectores y
regiones); crédito a empresas; empresas que nacen y cierran.\\
\textbf{Externo y fiscal}\newline\small(\emph{Cuentas externas y fiscales}) &
¿Cómo están las cuentas con el exterior y del Gobierno, y cómo se financian? &
Cuenta corriente por componentes; financiamiento (cuenta financiera); inversión extranjera directa y
remesas; deuda externa y posición de inversión internacional; ingresos, gastos, intereses, balance total y
primario, financiamiento y deuda del Gobierno.\\
\textbf{Comercio exterior}\newline\small(\emph{Qué vende y qué compra Colombia}) &
¿Cuánto, qué, para qué y con quién comercia Colombia? &
Exportaciones e importaciones en 12 meses; balanza comercial; composición por producto y por uso;
destinos y orígenes; precio frente a volumen; apertura y diversificación.\\
\bottomrule
\end{longtable}

La antigua dirección \texttt{/sectores/} redirige a la sección de sectores dentro de \emph{Crecimiento}.

\begin{leer}[Orden de lectura sugerido]
Para una primera lectura, siga el menú de izquierda a derecha: Ciclo $\rightarrow$ Crecimiento $\rightarrow$
Capacidad $\rightarrow$ Empleo $\rightarrow$ Inflación $\rightarrow$ Tasas $\rightarrow$ Curva TES
$\rightarrow$ Tasa de cambio $\rightarrow$ Sector empresarial $\rightarrow$ Externo y fiscal $\rightarrow$
Comercio exterior. Es la cadena causal que se desarrolla en la Parte~IV.
\end{leer}

% ---------------------------------------------------------------------
\chapter{La portada}
\label{es:u:portada}

La portada responde en una pantalla \emph{¿Cómo va la economía colombiana?} y sirve de índice para las
páginas de análisis. Tiene, de arriba abajo, siete bloques.

\section{Encabezado: dos fechas distintas}
Bajo el título aparecen dos fechas que conviene no confundir:
\begin{itemize}
\item \textbf{datos al}: la fecha del dato más reciente incorporado en el sitio (normalmente un dato de
  mercado diario);
\item \textbf{página actualizada}: el momento en que se construyó la página.
\end{itemize}
Que la página se haya actualizado hoy no significa que todas las cifras sean de hoy: cada cifra lleva su
propio período de referencia (capítulo~\ref{es:u:actualizacion}).

\section{El resumen en tres frases}
Un párrafo sin jerga, generado con reglas explícitas, dice: (1) si la economía crece por encima, al ritmo o por
debajo de su ritmo habitual; (2) si la inflación sube, baja o se mantiene, y dónde está frente a la meta; (3) en
qué nivel tiene el Banco de la República su tasa y qué busca con él (frenar la inflación, estimular la
economía o ninguna de las dos). Si la Junta Directiva ya anunció un cambio de tasa que aún no rige, la tercera
frase lo dice explícitamente. Las reglas exactas están en la sección~\ref{es:i:resumen}.

\section{El recuadro del ciclo}
Resume la página \emph{Ciclo} en cuatro cifras:
\begin{description}[style=nextline,leftmargin=0pt]
\item[Fase] Expansión, Desaceleración, Contracción o Recuperación, con una frase que la describe y el
  trimestre de referencia.
\item[Brecha] Cuánto produce la economía por encima ($+$) o por debajo ($-$) de su capacidad estimada, en
  porcentaje.
\item[Dirección] El cambio de la brecha en los dos últimos trimestres, en puntos porcentuales: positivo
  significa que la economía se aleja hacia arriba de su capacidad (o se acerca desde abajo).
\item[Racha] Cuántos trimestres seguidos lleva en la fase actual.
\end{description}
El enlace \emph{Ver el reloj del ciclo} abre la página de análisis.

\section{Las seis tarjetas principales}
Cada tarjeta tiene un título, una \emph{etiqueta de estado} de color, la cifra principal, su descripción y
período, una línea de comparación y un minigráfico de los dos últimos años. Toda la tarjeta enlaza a la página
de análisis correspondiente.

\begin{center}\small
\begin{tabularx}{\linewidth}{@{}lXX@{}}
\toprule
\textbf{Tarjeta} & \textbf{Cifra} & \textbf{Línea de comparación}\\\midrule
Crecimiento & PIB real, variación anual (trimestre) & Ritmo habitual (crecimiento potencial)\\
Inflación & Variación anual del IPC (mes) & Hace un año · meta 3\%\\
Tasa del Banco & Tasa de política vigente & Hace un año\\
Desempleo & Tasa de desempleo desestacionalizada (mes) & Hace un año\\
Dólar & TRM, pesos por dólar (día) & Variación en un año (\%)\\
Bolsa (COLCAP) & Puntos del índice (día) & Variación en un año (\%)\\
\bottomrule
\end{tabularx}
\end{center}

\section{El recuadro de la curva TES}
Un gráfico pequeño con la curva en pesos a 1, 5 y 10 años en la fecha más reciente y la misma curva un año antes,
junto con tres cifras: la \textbf{pendiente 10a $-$ 1a} (en puntos porcentuales), la \textbf{forma} de la curva
(\emph{Normal}, \emph{Plana} o \emph{Invertida}; reglas en la sección~\ref{es:i:estados-tabla}) y el
\textbf{cambio del TES a 10 años en 12 meses}. El enlace \emph{Explorar la curva día a día} abre la página
\emph{Curva TES}.

\section{``Lectura en un vistazo''}
Una tarjeta por tema (Ciclo, Crecimiento, Capacidad, Sectores, Empleo, Inflación, Tasas, Curva TES, Tasa de
cambio, Sector empresarial, Externo y fiscal, Comercio exterior). Cada una muestra:
\begin{itemize}
\item el grupo y el tema (por ejemplo \emph{Precios y tasas · Inflación});
\item cuando aplica, la \textbf{etiqueta de estado} del tema;
\item una frase que resume el tema con el dato más reciente;
\item \textbf{dos cifras clave}, cada una con su cambio (flecha \cmsym{▲} si sube, \cmsym{▼} si baja), la ventana
  del cambio (por ejemplo \emph{en 12 meses}) y su período de referencia;
\item el enlace \emph{Ver análisis}.
\end{itemize}

\subsection{Qué significa cada etiqueta de estado}
\label{es:u:pildoras}
Las etiquetas tienen tres tonos: \textcolor{cmgreen}{\textbf{verde}} (situación favorable o normal),
\textcolor{cmamber}{\textbf{ámbar}} (atención) y \textcolor{cmred}{\textbf{rojo}} (situación desfavorable);
las del dólar y la bolsa son neutras (gris), porque un peso más fuerte o una bolsa al alza no son, en sí mismos,
buenos o malos para todos los lectores.

\begin{longtable}{@{}L{0.17\linewidth}L{0.25\linewidth}L{0.50\linewidth}@{}}
\toprule
\textbf{Tema} & \textbf{Etiqueta} & \textbf{Regla (exacta)}\\\midrule\endhead
Crecimiento & Por encima de lo habitual & PIB anual $\ge$ crecimiento potencial $+1$\pp\ (verde)\\
 & Ritmo normal & entre potencial $-1$\pp\ y potencial $+1$\pp\ (verde)\\
 & Lento & PIB anual $\ge 0$ pero $<$ potencial $-1$\pp\ (ámbar)\\
 & En contracción & PIB anual $<0$ (rojo)\\\midrule
Inflación & En la meta & inflación anual entre 2\% y 4\% (verde)\\
 & Bajo la meta & inflación $<2\%$ (ámbar)\\
 & Sobre la meta & inflación $>4\%$ y $\le 5\%$ (ámbar)\\
 & Alta & inflación $>5\%$ (rojo)\\\midrule
Tasa del Banco & Alta: frena la economía & tasa real ex ante $>$ techo de la neutral $+0{,}5$\pp\ (ámbar)\\
 & Neutral & tasa real ex ante dentro de la neutral $\pm0{,}5$\pp\ (verde)\\
 & Baja: estimula la economía & tasa real ex ante $<$ piso de la neutral $-0{,}5$\pp\ (ámbar)\\
 & Anunciada: sube / baja & la Junta anunció un cambio que la serie oficial aún no registra\\\midrule
Desempleo & Mejorando & desempleo $<$ el de hace 12 meses $-0{,}3$\pp\ (verde)\\
 & Estable & dentro de $\pm0{,}3$\pp\ del de hace 12 meses (verde)\\
 & Empeorando & desempleo $>$ el de hace 12 meses $+0{,}3$\pp\ (ámbar)\\\midrule
Dólar & Peso más fuerte & TRM cae más de 5\% en 12 meses\\
 & Estable & TRM cambia entre $-5\%$ y $+5\%$\\
 & Peso más débil & TRM sube más de 5\% en 12 meses\\\midrule
Bolsa & Al alza / Estable / A la baja & COLCAP $>+5\%$, entre $\pm5\%$, $<-5\%$ en 12 meses\\\midrule
Ciclo & Expansión, Recuperación & verde\\
 & Desaceleración & ámbar\\
 & Contracción & rojo\\
\bottomrule
\end{longtable}

\begin{alerta}
Las etiquetas comparan un dato con un umbral fijo; no son juicios sobre la economía ni señales de compra o
venta. Una etiqueta puede cambiar por una diferencia mínima alrededor del umbral: lea siempre la cifra, no solo
el color.
\end{alerta}

\section{``Todos los indicadores''}
\label{es:u:tabla}
Una tabla de 16 indicadores (también en la página \emph{Datos}) con cinco columnas: \emph{Indicador},
\emph{Último dato}, \emph{Período}, \emph{Cambio} y \emph{Frente a su historia}.

\begin{itemize}
\item \textbf{Cambio.} Para tasas y porcentajes, en \textbf{puntos porcentuales} (pp); para niveles (TRM,
  tasa de cambio real, COLCAP), en \textbf{porcentaje}. La ventana es de 12 meses, salvo para la inflación que
  espera el mercado (1 año y largo plazo) y el TES a 10 años, que se comparan con 3 meses antes porque son
  datos de mercado que se mueven rápido.
\item \textbf{Frente a su historia.} Una barra con el \textbf{percentil histórico} del último dato frente a
  todas sus observaciones desde 2010 (fórmula en la sección~\ref{es:f:percentil}), traducido a palabras.
\end{itemize}

\begin{center}\small
\begin{tabular}{@{}lll@{}}
\toprule
\textbf{Percentil $p$} & \textbf{Palabra} & \textbf{Lectura}\\\midrule
$p<10$ & Muy bajo & entre el 10\% de valores más bajos desde 2010\\
$10\le p<30$ & Bajo & \\
$30\le p\le 70$ & Normal & en la franja central de su historia\\
$70< p\le 90$ & Alto & \\
$p>90$ & Muy alto & entre el 10\% de valores más altos desde 2010\\
\bottomrule
\end{tabular}
\end{center}

\begin{center}\small
\begin{longtable}{@{}p{0.55\linewidth}ll@{}}
\toprule
\textbf{Indicador} & \textbf{Cambio en} & \textbf{Ventana}\\\midrule\endhead
Crecimiento de la economía (PIB real, anual) & pp & 12 meses\\
Producción frente a su capacidad (brecha) & pp & 12 meses\\
Actividad económica mensual (ISE, anual) & pp & 12 meses\\
Desempleo & pp & 12 meses\\
Inflación anual & pp & 12 meses\\
Inflación de fondo (sin alimentos ni regulados) & pp & 12 meses\\
Inflación que espera el mercado (próximo año) & pp & 3 meses\\
Inflación que espera el mercado (largo plazo, 5y5y) & pp & 3 meses\\
Tasa de interés del Banco de la República & pp & 12 meses\\
Tasa de interés real (descontada la inflación esperada) & pp & 12 meses\\
Bonos del Gobierno a 10 años (TES) & pp & 3 meses\\
Dólar (TRM, pesos por dólar) & \% & 12 meses\\
Tasa de cambio real (2010 = 100) & \% & 12 meses\\
Bolsa de Colombia (COLCAP) & \% & 12 meses\\
Cuenta corriente (\% del PIB) & pp & 12 meses\\
Deuda bruta del Gobierno (\% del PIB) & pp & 12 meses\\
\bottomrule
\end{longtable}
\end{center}

\begin{leer}
``Muy alto'' no significa ``malo'' ni ``Muy bajo'' significa ``bueno'': un percentil solo ubica el dato dentro de
su propia historia. Una inflación ``Muy baja'' puede ser una buena noticia; un COLCAP ``Muy alto'' solo dice
que el índice está cerca de sus máximos desde 2010. Los extremos son los datos que más se alejan de su
normalidad y merecen una lectura más cuidadosa en su página.
\end{leer}

\section{Fuentes oficiales y publicaciones}
La portada termina con la tabla \emph{Fuentes oficiales} (fuente, frecuencia, último dato y estado: \emph{Al
día}, \emph{Retrasado} o \emph{Pendiente}; véase el capítulo~\ref{es:u:actualizacion}), los enlaces a la
metodología y al manual, y la nota \emph{Últimas publicaciones oficiales}: tarjetas breves con el dato más
reciente de cada entidad (por ejemplo, la última decisión de tasa, la TRM, el COLCAP, el desempleo o la tasa del
TES a 10 años), la fecha de publicación u observación y un enlace a la página oficial. Cuando la fuente no trae
la fecha de publicación, la tarjeta dice \emph{dato más reciente}. Al pie hay enlaces fijos a las salas de prensa
del Banco de la República, el DANE, el Ministerio de Hacienda y la Superintendencia Financiera.

\begin{leer}[Regla de oro de la portada]
Lea siempre el \textbf{período de referencia} junto a la cifra. En una misma pantalla conviven datos diarios
(tasas, dólar, bolsa), mensuales (inflación, desempleo, ISE), trimestrales (PIB, cuenta corriente) y anuales
(deuda del Gobierno). El sitio nunca rellena una frecuencia con otra.
\end{leer}

% ---------------------------------------------------------------------
\chapter{Anatomía de una página de análisis}
\label{es:u:anatomia}

Todas las páginas de análisis comparten la misma estructura. Conocerla una vez basta para leer las once.

\section{Encabezado y selector de horizonte}
El encabezado tiene la ruta de navegación, la fecha \emph{datos al}, el título de la página y un párrafo que
dice qué se encontrará en ella. Debajo está el \textbf{selector de horizonte}:
\emph{Ver los gráficos de los últimos 3 años · 5 años · 10 años · Todo} (sección~\ref{es:u:horizonte}).

\section{Secciones numeradas}
El cuerpo se divide en secciones numeradas (01, 02, \ldots), cada una titulada con una pregunta (por ejemplo
\emph{¿Qué sectores impulsan la economía?}). Cada sección tiene, en este orden:

\begin{enumerate}
\item \textbf{``Lo clave''}: de dos a cuatro frases generadas con los datos más recientes, con las cifras
  resaltadas. Responden la pregunta de la sección antes de mirar cualquier gráfico.
\item \textbf{Los gráficos} de la sección (o una tabla).
\end{enumerate}

La \textbf{primera sección} de cada página es un panel de medidas (\emph{La economía en cinco medidas},
\emph{El crecimiento en ocho medidas}, \emph{Las cuentas externas y fiscales en doce medidas}\ldots): un
conjunto de mosaicos, cada uno con el nombre de la medida, su valor, una línea de contexto (período,
comparación, referencia) y, cuando aplica, un tono de color con la misma lógica de las etiquetas. Los
mosaicos con enlace llevan a la sección donde se analiza esa medida.

La \textbf{última sección} de cada página, \emph{Bases metodológicas y literatura}, lista las referencias
académicas e institucionales que respaldan los métodos de la página, cada una con una línea que dice para qué
se usa (todas están compiladas en el capítulo~\ref{es:r:referencias}).

\section{Anatomía de un gráfico}
\label{es:u:grafico}
Cada gráfico es un bloque con estas partes:

\begin{description}[style=nextline,leftmargin=0pt]
\item[Título] Describe qué se mide, a veces con el período entre paréntesis (por ejemplo
  \emph{Informalidad por sector (may--jul)}).
\item[Franja de cifras] Encima del gráfico, una cifra por serie principal: nombre, último valor, su fecha y
  su cambio con flecha (\cmsym{▲}/\cmsym{▼}) frente a un punto de comparación explícito (\emph{vs. hace 1 año},
  \emph{vs. trimestre anterior}, \emph{vs. hace 1 mes}). Las tasas cambian en pp y los niveles en \%.
\item[Área del gráfico] Una escala vertical; leyenda arriba o al lado. Las franjas grises verticales, cuando
  aparecen, marcan recesiones fechadas en la página \emph{Ciclo}.
\item[Panel de cambios] En muchos gráficos de series de tiempo, un segundo panel debajo muestra, fecha por
  fecha, el cambio de la serie principal (barras azules cuando sube, naranjas cuando baja). Comparte el eje de
  fechas con el panel superior.
\item[Línea de lectura ``?''] Debajo del gráfico, un párrafo precedido de un círculo con ``?'' explica
  \emph{qué mirar}: qué es cada línea o barra y qué significa que suba o baje.
\item[Pie de fuente] \emph{Fuente:} entidad y publicación de donde salen los datos.
\item[Tooltip de metodología] Junto a la fuente, un botón redondo ``?''. Al pasar el cursor (o tocarlo en el
  teléfono) aparece un recuadro \emph{Metodología.} con la fórmula, los parámetros, la frecuencia y las
  advertencias del cálculo. La línea de lectura explica \emph{qué} mirar; la metodología, \emph{cómo} se calculó.
\end{description}

\section{La lupa: cómo leer e interpretar cada gráfico}
\label{es:u:lupa}
En la esquina de cada gráfico hay un botón con el icono de una lupa. Al pulsarlo se abre la ventana
\emph{Lupa · cómo leer este gráfico}, que reúne en un solo lugar todo lo necesario para interpretarlo. Sus
partes, siempre en el mismo orden, son:

\begin{enumerate}
\item \textbf{Título}: el del gráfico.
\item \textbf{Qué nos cuenta}: la idea central del gráfico en una o dos frases.
\item \textbf{Lo que muestra hoy}: una copia de las frases de ``Lo clave'' de la sección, con las cifras
  vigentes.
\item \textbf{Cómo se lee}: ejes, colores, líneas de referencia y convenciones.
\item \textbf{Por qué importa}: la relevancia económica y para un inversionista.
\item \textbf{Cómo interpretarlo}: una lista de lecturas típicas (qué significa cada configuración del
  gráfico) y de cautelas.
\item \textbf{Fórmulas}: cada fórmula con su nombre, su expresión y, cuando hace falta, una nota.
\item \textbf{Metodología y fuente}: el mismo texto del tooltip de metodología y la fuente.
\end{enumerate}

La ventana se cierra con \cmsym{✕}, con la tecla \emph{Esc} o haciendo clic fuera de ella. El contenido de cada
lupa se reproduce en el atlas de gráficos de este Codex (Parte~III).

\section{Ampliar un gráfico}
El botón \cmexp\ (\emph{Ampliar gráfico}) lleva el gráfico al centro de la pantalla, con más altura y el fondo
atenuado, sin mover la página. El gráfico ampliado conserva todos sus controles (selectores, clics, cajas
emergentes). Se cierra con el mismo botón (que pasa a \cmsym{✕}), con \emph{Esc} o haciendo clic en el fondo.

\section{El selector de horizonte y el rebase}
\label{es:u:horizonte}
Los botones \emph{3 años}, \emph{5 años}, \emph{10 años} y \emph{Todo} recortan \textbf{todos} los gráficos de
series de tiempo de la página a la misma ventana, que termina en el último dato. Detalles que importan:

\begin{itemize}
\item La opción predeterminada es \emph{10 años}; el navegador recuerda la última elección y la aplica en
  todas las páginas.
\item La escala vertical se ajusta a los datos visibles; por eso, con \emph{Todo} o con 10 años, episodios
  extremos como 2020 se ven completos y comprimen el resto, y con 3 años la escala se abre.
\item Si una serie empieza después del inicio de la ventana (por ejemplo, la informalidad desde 2021), el
  gráfico arranca en su primer dato, sin espacio vacío.
\item Los gráficos que no son series de tiempo (cortes transversales por sector, ciudad o plazo, mapas,
  tablas, la curva TES de un día) no cambian con el horizonte.
\item Los gráficos \textbf{base 100} (por ejemplo, COLCAP frente al equiponderado y a las 7 Magníficas, o el
  peso frente a otras monedas) se \textbf{re-basan}: cada serie vale 100 en el primer dato visible de la ventana,
  \[
  \tilde I_t = 100\,\frac{I_t}{I_{t_0}},\qquad t_0=\text{primer dato dentro del horizonte},
  \]
  de modo que el valor de cada línea al final es directamente su rendimiento acumulado en la ventana. Cambiar
  el horizonte cambia los niveles de las líneas, pero no su orden relativo dentro de la ventana.
\end{itemize}

\begin{alerta}[El horizonte cambia la impresión visual]
Una misma serie puede verse ``alta'' con 3 años y ``normal'' con \emph{Todo}. Antes de sacar una conclusión
sobre niveles, revise el horizonte y, si es posible, compárelo con el percentil histórico de la tabla de
indicadores (que siempre usa la historia desde 2010).
\end{alerta}

\section{Cajas emergentes (hover)}
Al pasar el cursor por un punto, barra o casilla aparece una caja con su fecha, su valor y, en muchos gráficos,
su cambio frente al mismo punto de comparación de la franja (por ejemplo \emph{\cmsym{▲} +0,5 pp vs. trimestre
anterior}). Así se puede leer el cambio en cualquier fecha histórica, no solo en el último dato. En los gráficos
con dos paneles, la caja muestra el valor del panel superior y el cambio del inferior en la misma fecha. En los
mapas de calor, la caja muestra el valor exacto aunque la escala de color esté recortada.

\section{Controles interactivos propios de algunas páginas}
\begin{itemize}
\item \textbf{Reloj del ciclo} (\emph{Ciclo}): una barra para elegir el trimestre, botones de recorrido
  (\emph{1 año}, \emph{2 años}, \emph{3 años}) y \emph{Volver al último dato}; al pasar el cursor por el gráfico de
  barras del recorrido, el reloj se mueve a ese trimestre.
\item \textbf{Selector de trimestre de los sectores} (\emph{Crecimiento}): \emph{Ver el trimestre:} muestra
  cualquier trimestre desde 2010; un clic en el mapa de calor lleva las barras a esa columna.
\item \textbf{Explorador de la curva TES} (\emph{Curva TES}): interruptor pesos/UVR, barra de fechas desde 2003,
  años para comparar (\emph{Comparar con el cierre de:}) y la historia de 1, 5 y 10 años (pasar el cursor mueve
  la curva; un clic fija la fecha).
\item \textbf{Mapas regionales} (\emph{Capacidad}) y \textbf{listas desplegables} de ciudades (\emph{Empleo}).
\end{itemize}

\section{La ventana ``Para entender''}
Varias páginas tienen un botón que abre una ventana explicativa sobre la fuente y los conceptos de la página:
cómo se miden exportaciones e importaciones (FOB, CIF, CUODE); qué es la curva cero cupón; de dónde salen las cifras
de las empresas; cómo leer la balanza de pagos y las cuentas del Gobierno; qué son las 188 subclases del IPC;
cómo se calcula la TRM y qué es la tasa de cambio real; cómo funciona la tasa del Banco y cómo llega a la
economía. Cada ventana termina con enlaces a las fuentes oficiales.

\section{Idioma, tema y teléfono}
\begin{itemize}
\item \textbf{Idioma.} El enlace \emph{English}/\emph{Español} cambia a la misma página en el otro idioma.
  Cambian los textos y el formato numérico: coma decimal y punto de miles en español (3,5\%; 3.313), punto
  decimal y coma de miles en inglés (3.5\%; 3,313). Los trimestres se escriben T1--T4 en español y Q1--Q4 en
  inglés.
\item \textbf{Tema oscuro.} El botón de tema cambia entre claro y oscuro; los colores de los gráficos se
  traducen a su par en el otro tema, de modo que el significado de cada color (azul sube, naranja baja; verde,
  ámbar y rojo en las etiquetas) se mantiene. El navegador recuerda la elección.
\item \textbf{Teléfono.} Tarjetas y gráficos se apilan en una columna; el menú se abre con \cmsym{☰}; los
  tooltips de metodología se abren tocando el ``?''.
\end{itemize}

% ---------------------------------------------------------------------
\chapter{Frecuencia de actualización y estado de las fuentes}
\label{es:u:actualizacion}

\section{Cuándo se actualiza el sitio}
Los datos se descargan automáticamente de las fuentes oficiales: las series diarias (tasas, dólar, bolsa), cada
hora en días hábiles; todas las fuentes, cada noche. Antes de publicarse, cada tabla pasa controles de
calidad (identidad de la serie, rangos plausibles, fechas sin huecos ni duplicados, recálculo de fórmulas). Si
una fuente falla, se conserva su último dato válido y su estado lo refleja.

\section{Período medido frente a fecha de publicación}
La fecha que acompaña a cada cifra es el \textbf{período medido} (el trimestre del PIB, el mes del IPC, el día
de la TRM), no el día en que se publicó. Por eso, un PIB fechado en el segundo trimestre puede ser el dato más
reciente disponible varios meses después de terminado ese trimestre.

\section{Rezagos típicos de publicación}
\begin{center}\small
\begin{tabularx}{\linewidth}{@{}llX@{}}
\toprule
\textbf{Variable (fuente)} & \textbf{Frecuencia} & \textbf{Rezago aproximado de publicación}\\\midrule
PIB y PIB por sectores (DANE) & Trimestral & unos 45 días después del trimestre\\
ISE (DANE) & Mensual & unos 45--50 días después del mes\\
IPC (DANE) & Mensual & primera semana del mes siguiente\\
Desempleo, GEIH (DANE) & Mensual & alrededor de un mes\\
Informalidad, GEIH (DANE) & Trimestre móvil & alrededor de mes y medio\\
Comercio exterior (DANE--DIAN) & Mensual & entre uno y dos meses\\
Tasa de política, TRM, TES, COLCAP, IBR & Diaria & el mismo día o el siguiente\\
Tasas de crédito y CDT (BanRep) & Semanal & pocos días\\
Balanza de pagos, deuda externa (BanRep) & Trimestral & unos dos a tres meses\\
Deuda del Gobierno (MinHacienda vía BanRep) & Anual & varios meses después del cierre\\
10.000 empresas (Supersociedades) & Anual & a mediados del año siguiente\\
\bottomrule
\end{tabularx}
\end{center}

\section{Qué significa ``Al día''}
\label{es:u:aldia}
La tabla de fuentes asigna a cada una uno de tres estados:
\begin{description}[style=nextline,leftmargin=0pt]
\item[Al día] El último período disponible está dentro de la tolerancia esperada para su frecuencia: han
  pasado menos días que el umbral de la tabla siguiente desde la fecha con la que se registra ese período.
\item[Retrasado] Han pasado más días que el umbral: la fuente no ha publicado un dato nuevo cuando se esperaba,
  o la descarga falló y se conserva el último dato válido.
\item[Pendiente] La fuente aún no se ha descargado por primera vez.
\end{description}

\begin{center}\small
\begin{tabular}{@{}lr@{}}
\toprule
\textbf{Fuente} & \textbf{Umbral (días)}\\\midrule
Bolsa -- COLCAP (diaria) & 7\\
Tasas de los TES, tasa de política, TRM (diarias) & 10\\
Canasta COLCAP y acciones del COLCAP & 10\\
Inflación -- IPC (mensual) & 45\\
Inflación básica (mensual) & 75\\
Desempleo -- GEIH (mensual) & 100\\
Tasa de cambio real -- ITCR (mensual) & 100\\
Actividad mensual -- ISE; Informalidad -- GEIH & 110\\
PIB real (trimestral) & 130\\
Cuenta corriente (trimestral) & 200\\
PIB por sectores (trimestral) & 230\\
Deuda del Gobierno (anual) & 800\\
\bottomrule
\end{tabular}
\end{center}

\begin{alerta}[``Al día'' no es ``de hoy'']
\emph{Al día} significa que la fuente cumple su calendario normal, no que el dato sea reciente en sentido
absoluto. Un PIB \emph{Al día} puede referirse a un trimestre que terminó hace más de tres meses, y una deuda
anual \emph{Al día} puede referirse al cierre del año anterior. Los umbrales son amplios porque se cuentan
desde la fecha con la que se registra el período (normalmente su primer día), no desde su publicación.
\end{alerta}
